% ============================================================
% GENERAL CONCLUSION (Following Tarek's Conclusion Générale)
% ============================================================

\chapter*{General Conclusion}
\addcontentsline{toc}{chapter}{General Conclusion}
\markboth{General Conclusion}{General Conclusion}

% ============================================================
\section*{Project Summary}
\addcontentsline{toc}{section}{Project Summary}

At the end of this end-of-studies project, we successfully designed and developed a DevSecOps Policy-as-Code microservice dedicated to the automated security analysis of infrastructure code prior to deployment. Through four consecutive sprints, we progressively built a coherent and performant security scanning ecosystem, combining automated analysis, educational feedback, and CI/CD integration.

The iterative approach adopted from the beginning of the project proved particularly well-suited, enabling modular evolution and fluid integration of new functionalities. The 4-layer architecture facilitated the parallel development of different system components, optimizing our development cycle while maintaining a high level of quality and coherence.

% ============================================================
\section*{Summary of Achievements}
\addcontentsline{toc}{section}{Summary of Achievements}

\subsection*{Sprint 1: Core Microservice Foundation (Feb 3--14, 2026)}

During the first sprint, we established the solid foundations of the project by:

\begin{itemize}
    \item Conducting a comprehensive analysis of the DevSecOps landscape and existing tools (Checkov, tfsec, Trivy, OPA).
    \item Defining the 4-layer architecture (API, Parser, Normalizer, Policy Engine).
    \item Setting up the Python 3.10 development environment with FastAPI.
    \item Creating the basic API with the \texttt{POST /analyze} endpoint responding to requests.
    \item Designing the initial set of 12 core security policies based on real-world breach analysis.
    \item Configuring the Git repository and development workflow.
\end{itemize}

\noindent\textbf{Deliverables:} Functional FastAPI service with basic endpoint, architecture document, list of policies with justifications.

% TODO: INSERT SCREENSHOT OF BASIC API RESPONDING (Swagger /docs page)
% \begin{figure}[H]
%     \centering
%     \includegraphics[width=0.8\textwidth]{figures/sprint1-api.png}
%     \caption{Sprint 1 — FastAPI Swagger documentation}
%     \label{fig:sprint1-api}
% \end{figure}

\subsection*{Sprint 2: Policy Engine and 100 Security Rules (Feb 17--28, 2026)}

The second sprint enriched the platform with its core parsing and normalization capabilities:

\begin{itemize}
    \item Implemented the \textbf{TerraformParser} using \texttt{python-hcl2}, capable of parsing HCL syntax and extracting resources with their types, names, and attributes.
    \item Implemented the \textbf{DockerfileParser} with custom regex-based parsing, processing each instruction line by line with line number tracking.
    \item Implemented the \textbf{KubernetesParser} using \texttt{PyYAML}, extracting pods, containers, and security contexts from Deployment and Pod specifications.
    \item Designed and implemented the \textbf{Normalizer} class that converts all three heterogeneous formats into a unified \texttt{Resource} model—the key architectural innovation of the project.
    \item Created the \texttt{NormalizedCode} and \texttt{Resource} data models.
\end{itemize}

The normalizer design was the most significant architectural challenge of the project. Initially, we attempted to make policies format-aware, which would have resulted in 300 separate implementations (100 policies $\times$ 3 formats). After three days of design iteration, we arrived at the normalizer pattern, reducing this to 100 implementations—a \textbf{67\% reduction in code complexity}.

\noindent\textbf{Deliverables:} Three working parsers, unified normalizer, Resource data model.

% TODO: INSERT SCREENSHOT OF CODE - Normalizer converting Terraform to Resource model
% \begin{figure}[H]
%     \centering
%     \includegraphics[width=0.8\textwidth]{figures/sprint2-normalizer.png}
%     \caption{Sprint 2 — Normalizer converting Terraform code to unified Resource model}
%     \label{fig:sprint2-normalizer}
% \end{figure}

\subsection*{Sprint 3: Dashboard, Reports and Scan History (Mar 2--13, 2026)}

The third sprint was the most intensive, involving the implementation of the complete policy engine and all 100 security rules:

\begin{itemize}
    \item Created the \textbf{BasePolicy} abstract class with the \texttt{check()} method, establishing the Template Method design pattern.
    \item Implemented the \textbf{PolicyEngine} orchestrator that runs all policies, aggregates violations, and produces ALLOW/BLOCK decisions.
    \item Implemented all \textbf{100 security policies} organized across 14 categories:
    \begin{itemize}
        \item 12 core policies (hardcoded secrets, public S3/database, unencrypted storage, etc.)
        \item 88 extended policies covering AWS S3, EC2, RDS, IAM, VPC, Docker, Kubernetes, CloudTrail, ELB/WAF, Terraform, EKS/Lambda/ECS, messaging/cache, ECR/CloudFront.
    \end{itemize}
    \item Added detailed violation messages with severity levels, risk explanations, and remediation recommendations for every policy.
    \item Implemented the configurable severity threshold system (CRITICAL, HIGH, MEDIUM, LOW).
    \item Added the diff-only scanning feature using \texttt{parse\_diff\_lines()} to filter violations to only changed lines.
\end{itemize}

\noindent\textbf{Deliverables:} PolicyEngine with 100 working policies, configurable thresholds, diff scanning.

% TODO: INSERT SCREENSHOT OF SCAN RESULT showing violations found
% \begin{figure}[H]
%     \centering
%     \includegraphics[width=0.8\textwidth]{figures/sprint3-violations.png}
%     \caption{Sprint 3 — Policy engine scan result showing violations}
%     \label{fig:sprint3-violations}
% \end{figure}

\subsection*{Sprint 4: CI/CD Integration, Docker Deployment and Demo (Mar 16 -- Apr 10, 2026)}

The final and longest sprint focused on integration, visualization, and demonstration:

\begin{itemize}
    \item Created the \textbf{Jenkinsfile} for automated pipeline integration, with stages for checkout, security scan, build, and deploy.
    \item Set up \textbf{Jenkins} on localhost:8080 for CI/CD automation.
    \item Set up \textbf{Gitea} on localhost:3000 as a local Git server for offline demonstration.
    \item Implemented the \textbf{SQLite scan history} system with \texttt{save\_scan()}, \texttt{get\_history()}, and \texttt{get\_stats()} functions.
    \item Developed the \texttt{/metrics} endpoint tracking total requests, blocks, allows, block rate, top violations, and requests by code type.
    \item Built the \textbf{real-time dashboard} (HTML/CSS/JavaScript) served at \texttt{GET /}, featuring:
    \begin{itemize}
        \item Live statistics cards (total scans, blocked, allowed, block rate)
        \item Ring chart visualization of block/allow ratio
        \item Top violations ranking
        \item Scan history table with timestamps
    \end{itemize}
    \item Implemented the \textbf{HTML security report generator} with a dark-themed, professional design including severity badges, violation cards, and fix recommendations.
    \item Created the Docker containerization (\texttt{Dockerfile} + \texttt{docker-compose.yml}).
    \item Developed CI/CD integration examples for Jenkins, GitLab CI, and GitHub Actions.
    \item Built the \textbf{multi-developer offline demonstration} (Alice and Bob scenario).
\end{itemize}

% TODO: INSERT SCREENSHOTS:
% 1. Dashboard UI
% 2. HTML security report
% 3. Jenkins pipeline (green/red stages)
% 4. Gitea repository view
% 5. API Swagger docs

% \begin{figure}[H]
%     \centering
%     \includegraphics[width=0.9\textwidth]{figures/sprint4-dashboard.png}
%     \caption{Sprint 4 — Real-time security dashboard}
%     \label{fig:sprint4-dashboard}
% \end{figure}

% \begin{figure}[H]
%     \centering
%     \includegraphics[width=0.9\textwidth]{figures/sprint4-report.png}
%     \caption{Sprint 4 — HTML security report with violation details}
%     \label{fig:sprint4-report}
% \end{figure}

% \begin{figure}[H]
%     \centering
%     \includegraphics[width=0.9\textwidth]{figures/sprint4-jenkins.png}
%     \caption{Sprint 4 — Jenkins pipeline with security scan stages}
%     \label{fig:sprint4-jenkins}
% \end{figure}

\noindent\textbf{Demo scenario results:}
\begin{itemize}
    \item \textbf{Alice's insecure code}: 19 violations detected $\rightarrow$ \textcolor{red}{\textbf{BLOCK}} — deployment prevented.
    \item \textbf{Bob's secure code}: 0 violations detected $\rightarrow$ \textcolor{green!50!black}{\textbf{ALLOW}} — deployment proceeds.
    \item Dashboard shows: 2 scans, 1 blocked, 1 allowed, 50\% block rate.
\end{itemize}

\noindent\textbf{Deliverables:} Complete CI/CD integration, real-time dashboard, HTML reports, SQLite history, Docker deployment, offline demo.

% ============================================================
\section*{Technical Architecture}
\addcontentsline{toc}{section}{Technical Architecture}

The technical architecture deployed for our service stands out through its modularity and robustness. The 4-layer architecture encapsulates each major responsibility in an independent and specialized layer, communicating through standardized interfaces.

The \textbf{Normalizer} constitutes the architectural keystone of the system. By converting three disparate formats (HCL, Dockerfile, YAML) into a common Resource model, policy complexity was reduced by 67\% while maintaining multi-format support. This pattern could be applied to other multi-format analysis tools beyond this project.

The \textbf{Strategy pattern} enables seamless addition of new parsers, while the \textbf{Template Method pattern} in the BasePolicy class allows new security policies to be added by implementing a single \texttt{check()} method. The \textbf{Factory pattern} in \texttt{get\_parser()} routes to the appropriate parser based on code type.

The SQLite-backed history system provides lightweight, serverless persistence that survives service restarts without requiring external database infrastructure.

% ============================================================
\section*{Performance Results}
\addcontentsline{toc}{section}{Performance Results}

\begin{table}[H]
\centering
\caption{Performance benchmarks}
\label{tab:performance}
\begin{tabular}{|l|c|c|}
\hline
\textbf{Metric} & \textbf{Our Service} & \textbf{Target} \\
\hline
Average response time & 46ms & $<$50ms \cmark \\
\hline
Memory footprint & $\sim$150 MB & $<$200 MB \cmark \\
\hline
Startup time & $<$2 seconds & $<$5 seconds \cmark \\
\hline
Security policies & 100 & 100 \cmark \\
\hline
IaC formats supported & 3 & 3 \cmark \\
\hline
Offline capable & Yes & Yes \cmark \\
\hline
\end{tabular}
\end{table}

\begin{table}[H]
\centering
\caption{Speed comparison with existing tools}
\label{tab:speed-comparison}
\begin{tabular}{|l|c|c|}
\hline
\textbf{Tool} & \textbf{Response Time} & \textbf{Comparison} \\
\hline
\textbf{Our Service} & \textbf{46ms} & — \\
\hline
Checkov & 150--300ms & 3--6$\times$ slower \\
\hline
tfsec & $\sim$150ms & $\sim$3$\times$ slower \\
\hline
Trivy & $\sim$200ms & $\sim$4$\times$ slower \\
\hline
\end{tabular}
\end{table}

% ============================================================
\section*{Challenges Overcome}
\addcontentsline{toc}{section}{Challenges Overcome}

The development of a project of this scope required overcoming several major challenges:

\begin{enumerate}
    \item \textbf{Normalizer Design (3 days)}: The most significant challenge was designing the normalizer layer. Initial attempts to make policies format-aware resulted in explosively complex code. After three days of design iteration, the unified Resource model emerged, saving weeks of policy implementation effort.
    
    \item \textbf{Parser Complexity}: Terraform's HCL syntax proved too complex for custom parsing. We switched to the \texttt{python-hcl2} library, leveraging expert-maintained code and focusing our innovation on policy logic rather than parsing.
    
    \item \textbf{False Positive Management}: The \texttt{EXPENSIVE\_INSTANCE\_TYPE} policy initially flagged too many legitimate use cases. We addressed this by making policies configurable, adjusting severity levels, and adding contextual messages.
    
    \item \textbf{Scaling to 100 Policies}: The Template Method pattern in BasePolicy made it possible to implement 100 policies efficiently. Each policy is a self-contained class with a single \texttt{check()} method, making the codebase maintainable despite its size ($\sim$5,900 lines).
\end{enumerate}

% ============================================================
\section*{Future Perspectives}
\addcontentsline{toc}{section}{Future Perspectives}

While this project has achieved an advanced level of maturity, several evolution paths are envisioned:

\textbf{Short-term (1--2 months):}
\begin{itemize}
    \item Add Azure and GCP security policies (estimated: 1--2 days per cloud provider).
    \item Improve secret detection with entropy analysis and context awareness.
    \item Develop a VS Code extension for inline scanning while coding.
\end{itemize}

\textbf{Medium-term (3--6 months):}
\begin{itemize}
    \item Integrate machine learning for anomaly detection and intelligent policy suggestions.
    \item Add Grafana integration for advanced metrics visualization.
    \item Implement Slack/Teams notifications for scan results.
    \item Support additional IaC formats (Pulumi, AWS CDK).
\end{itemize}

\textbf{Long-term (6--12 months):}
\begin{itemize}
    \item Develop a SaaS offering (cloud-hosted version).
    \item Add enterprise features: role-based access control, compliance reports, audit logs.
    \item Implement AI-powered auto-fix suggestions.
\end{itemize}

% ============================================================
\section*{Personal and Professional Contribution}
\addcontentsline{toc}{section}{Personal and Professional Contribution}

This end-of-studies project represented an extremely enriching experience, both technically and professionally. It enabled us to deepen our skills in:

\begin{itemize}
    \item \textbf{Technical mastery}: Python advanced features (async, abstract classes, decorators), API design with FastAPI, testing with pytest, Docker containerization, and CI/CD pipeline engineering with Jenkins.
    \item \textbf{Security thinking}: Understanding threat models, analyzing real-world breach patterns, and translating security requirements into automated policy checks.
    \item \textbf{Architecture design}: The experience of designing the normalizer taught the value of investing time in architecture upfront—three days of design work saved weeks of implementation effort.
    \item \textbf{Project management}: Applying Scrum methodology in a solo context reinforced self-discipline, iterative planning, and scope management.
    \item \textbf{Communication}: Learning to explain complex technical concepts clearly through educational violation messages.
\end{itemize}

% ============================================================
\section*{Final Word}
\addcontentsline{toc}{section}{Final Word}

At the conclusion of this project, we are proud to have successfully met the challenge of designing and developing a functional, performant, and educational DevSecOps microservice. The four consecutive sprints produced a product that fully meets the initial objectives while incorporating advanced features such as 100 security policies, a real-time dashboard, and a complete offline CI/CD demonstration.

The 4-layer architecture adopted throughout the development proved its relevance, offering the flexibility necessary to evolve each component of the system independently. The normalizer design, in particular, stands as the project's most significant architectural contribution—demonstrating that good architecture decisions compound over time.

This project constitutes not only the culmination of our academic training but also a genuine springboard toward our professional future. Security does not have to slow teams down. With the right tools, we can deploy faster \textit{and} more securely. This project proves that automation, education, and speed can coexist in DevSecOps.
